\documentclass[12pt,a4paper]{article}
\usepackage[utf8]{inputenc}
\usepackage[T1]{fontenc}
\usepackage{geometry}
\usepackage{booktabs}
\usepackage{amsmath}
\usepackage{graphicx}
\usepackage{xcolor}
\usepackage{array}
\usepackage{longtable}
\usepackage{multirow}
\usepackage{hyperref}
\usepackage{fancyhdr}
\usepackage{titlesec}
\usepackage{enumitem}
\usepackage{caption}
\geometry{margin=2.5cm}
\definecolor{airblue}{RGB}{30,100,180}
\definecolor{evapgreen}{RGB}{34,139,34}
\definecolor{chilledred}{RGB}{180,30,30}
\definecolor{lightgray}{RGB}{245,245,245}
\pagestyle{fancy}
\fancyhf{}
\rhead{FYP: Data Center Cooling Strategies}
\lhead{Aleeza Hussain}
\cfoot{\thepage}
\renewcommand{\headrulewidth}{0.4pt}
\hypersetup{colorlinks=true,linkcolor=black,urlcolor=blue,citecolor=black}
\begin{document}
\begin{titlepage}
\centering
\vspace*{2cm}
{\LARGE\bfseries FYP: Comparative Analysis of Data Center Cooling Strategies\\[0.5em]}
{\large\bfseries Simulation Results \& Analysis Report\\[2em]}
{\large Aleeza Hussain\\[0.5em]}
{\normalsize April 2026\\[3em]}
\rule{\textwidth}{1pt}\\[1em]
{\normalsize\itshape
This report presents the full simulation results from three high-fidelity physics-based cooling models:\\
Air-Side Economization, Evaporative Cooling, and Chilled Water Systems.\\
All simulations run over 8{,}760 hours (one full year) using CloudSim Plus 8.x co-simulation.}\\[1em]
\rule{\textwidth}{1pt}
\vfill
{\small Department of Computer Science \\ Final Year Project --- April 2026}
\end{titlepage}
\tableofcontents
\newpage

%% ============================================================
\section{Executive Summary}
%% ============================================================

This report presents the complete simulation results for three data center cooling strategies evaluated under identical AI-era workload conditions over a full 8{,}760-hour annual cycle. The simulation framework integrates CloudSim Plus 8.x Discrete Event Simulation with custom-built physics engines implemented in Java 17 / Spring Boot 3.x, as described in the accompanying Methodology document.

\subsection*{Key Findings at a Glance}

\begin{table}[h!]
\centering
\caption{Annual Performance Summary --- All Three Cooling Techniques}
\label{tab:exec_summary}
\renewcommand{\arraystretch}{1.4}
\begin{tabular}{lrrrrr}
\toprule
\textbf{Metric} & \textbf{Air-Side Econ.} & \textbf{Evaporative} & \textbf{Chilled Water} & \textbf{Best} \\
\midrule
Total Energy (kWh/yr) & 478{,}090 & 141{,}539 & 288{,}475 & Evaporative \\
IT Energy (kWh/yr) & 372{,}125 & 141{,}173 & 248{,}000 & Evaporative \\
Average PUE & 1.285 & 1.002 & 1.163 & Evaporative \\
Average CUE (kgCO$_2$/kWh$_{IT}$) & 0.591 & 0.461 & 0.549 & Evaporative \\
Annual OpEx (USD) & \$97{,}033 & \$20{,}523 & \$42{,}553 & Evaporative \\
Carbon Emissions (kg/yr) & 219{,}921 & 65{,}108 & 158{,}196 & Evaporative \\
Water Usage (L/yr) & 0 & 0 & 515{,}331 & Air-Side \\
CAPEX (USD) & \$45{,}000 & N/A & \$630{,}000 & Air-Side \\
Payback Period (yrs) & 491.9 & N/A & 50.0 & --- \\
\bottomrule
\end{tabular}
\end{table}

\noindent\textbf{Overall Recommendation:} The ML-based recommender system, trained on the simulation outputs, recommends \textbf{Evaporative Cooling} as the optimal strategy for the simulated conditions, achieving the lowest PUE (1.002), lowest annual OpEx (\$20{,}523), and lowest carbon emissions (65{,}108 kg CO$_2$/yr). However, it carries a critical thermal compliance failure --- rack inlet temperatures reached 38.5\textdegree C, exceeding the ASHRAE Class A1 limit of 27\textdegree C --- making it unsuitable as a standalone solution for high-density AI workloads without supplemental mechanical cooling.

\newpage

%% ============================================================
\section{Simulation Architecture \& Co-Simulation Design}
%% ============================================================

\subsection{System Architecture Overview}

The simulation platform is built on a lock-step co-simulation model that tightly couples a Discrete Event Simulation (DES) layer with a physics-based thermal engine. Figure~\ref{fig:arch} describes the data flow.

\begin{table}[h!]
\centering
\caption{Co-Simulation Layer Responsibilities}
\label{tab:arch_layers}
\renewcommand{\arraystretch}{1.3}
\begin{tabular}{p{3.5cm}p{5cm}p{5cm}}
\toprule
\textbf{Layer} & \textbf{Technology} & \textbf{Responsibility} \\
\midrule
Computational Layer & CloudSim Plus 8.x & VM provisioning, task scheduling, hourly $P_{IT}$ output \\
Thermal Layer & Spring Boot 3.x / Java 17 & Physics engine: airflow, heat balance, COP calculations \\
Synchronisation & Lock-step clock ($\Delta t = 3600$s) & Pulls $P_{IT}$ each tick; returns $P_{cooling}$ to global tracker \\
Economic Layer & ProjectionEngine.java & 5-year NPV, carbon tax escalation, TOU tariffs \\
ML Recommender & Random Forest (scikit-learn) & Ranks techniques by composite score \\
\bottomrule
\end{tabular}
\end{table}

\subsection{Workload Model}

All three simulations use an identical AI-aware composite workload:
\[
U_{total}(t) = \bigl[U_{diurnal}(t) \times M_{AI}\bigr] + \sigma(t)
\]
where $U_{diurnal}$ is a sinusoidal function peaking at 14:00, $M_{AI} = 1.3$ (Mixed AI workload), and $\sigma(t)$ is Gaussian noise ($\pm 5\%$). Weekend load is reduced by 70\%. UPS and PDU losses are modelled with quadratic efficiency curves.

\subsection{Simulation Parameters}

\begin{table}[h!]
\centering
\caption{Common Simulation Input Parameters}
\label{tab:sim_params}
\renewcommand{\arraystretch}{1.3}
\begin{tabular}{llll}
\toprule
\textbf{Category} & \textbf{Parameter} & \textbf{Value} & \textbf{Source} \\
\midrule
Thermal & ASHRAE Max Inlet ($T_{max}$) & 27.0\textdegree C & ASHRAE TC 9.9 \\
Thermal & Thermal Mass (Rack) & 15.0 kJ/K & Design Default \\
Psychrometric & Air Density ($\rho$) & 1.2 kg/m$^3$ & Ideal Gas Law \\
Psychrometric & Latent Heat ($L_v$) & 2{,}260 kJ/kg & Water Vaporisation \\
Power & Chiller Reference COP & 3.0 & Conservative Baseline \\
Economic & Base Electricity Rate & \$0.12/kWh & Regional Default \\
Carbon & Carbon Tax (2030) & \$254/ton CO$_2$ & IPCC Pathway \\
Workload & AI Training Multiplier & 1.8$\times$ & Sustained High Load \\
Control & Mode Switch Hysteresis & 2.0\textdegree C & Stability Buffer \\
Simulation & Total Hours & 8{,}760 & Full Annual Cycle \\
\bottomrule
\end{tabular}
\end{table}

\newpage

%% ============================================================
\section{Air-Side Economizer: Simulation Results}
%% ============================================================

\subsection{System Description}

The Air-Side Economizer leverages ambient outdoor air to offset mechanical cooling. A decision-tree state machine selects between three operating modes based on outdoor dry-bulb temperature ($T_{out}$) and relative humidity ($RH_{out}$):

\begin{itemize}[noitemsep]
  \item \textbf{FULL\_ECON}: $T_{out} \leq 24$\textdegree C \textit{and} $RH_{out} \leq 60\%$ --- 100\% outdoor air, compressor off.
  \item \textbf{PARTIAL\_TRIM}: $T_{out} \leq 24$\textdegree C \textit{but} $RH_{out} > 60\%$ --- 20\% outdoor air, partial mechanical assist.
  \item \textbf{MECHANICAL\_ONLY}: $T_{out} > 24$\textdegree C --- 0\% outdoor air, full compressor operation.
\end{itemize}

A 15\% filter penalty is applied to fan power in economiser modes (MERV filtration resistance). Required volumetric airflow is:
\[
V_{req}\ (\text{CFM}) = \frac{P_{IT,kW} \times 3160}{\rho \cdot C_p \cdot \Delta T}
\]

\subsection{Annual Energy Results (8{,}760-Hour Simulation)}

\begin{table}[h!]
\centering
\caption{Air-Side Economizer --- Annual Energy \& Emissions Summary}
\label{tab:air_annual}
\renewcommand{\arraystretch}{1.3}
\begin{tabular}{lrr}
\toprule
\textbf{Metric} & \textbf{Value} & \textbf{Unit} \\
\midrule
Total IT Energy & 372{,}125 & kWh \\
Total Cooling Energy & 105{,}965 & kWh \\
Total Facility Energy & 478{,}090 & kWh \\
Average PUE & 1.285 & --- \\
Average CUE & 0.591 & kgCO$_2$/kWh$_{IT}$ \\
Total Carbon Emissions & 219{,}921 & kg CO$_2$ \\
Carbon Savings vs. Baseline & 88{,}198 & kg CO$_2$ \\
Energy Savings vs. Baseline & 28.6\% & --- \\
Water Usage & 0 & Litres \\
\bottomrule
\end{tabular}
\end{table}

\subsection{Economic Results}

\begin{table}[h!]
\centering
\caption{Air-Side Economizer --- Economic Analysis}
\label{tab:air_econ}
\renewcommand{\arraystretch}{1.3}
\begin{tabular}{lrr}
\toprule
\textbf{Metric} & \textbf{Value} & \textbf{Unit} \\
\midrule
CAPEX & \$45{,}000 & USD \\
Annual Electricity Cost & \$69{,}323 & USD \\
Annual Carbon Tax Cost & \$27{,}710 & USD \\
Total Annual OpEx & \$97{,}033 & USD \\
Annual Savings vs. Baseline & \$91 & USD \\
Simple Payback Period & 491.9 & Years \\
\bottomrule
\end{tabular}
\end{table}

\noindent\textbf{Note on Payback Period:} The extremely long payback period (491.9 years) reflects the high carbon tax burden (\$27{,}710/yr) applied at the IPCC 2030 rate of \$254/ton CO$_2$. The system's carbon emissions (219{,}921 kg/yr) are high because the simulation site uses a carbon-intensive grid (0.55 kgCO$_2$/kWh). The electricity cost savings are modest because the system still requires significant mechanical cooling during hot/humid hours.

\subsection{5-Year Financial Projection}

\begin{table}[h!]
\centering
\caption{Air-Side Economizer --- 5-Year Cost Projection (with 15\% Carbon Tax Escalation)}
\label{tab:air_5yr}
\renewcommand{\arraystretch}{1.3}
\begin{tabular}{rrrrrr}
\toprule
\textbf{Year} & \textbf{Energy (kWh)} & \textbf{Elec. Cost} & \textbf{Carbon Tax} & \textbf{Total Cost} & \textbf{Cum. Savings} \\
\midrule
1 & 435{,}449 & \$65{,}317 & \$3{,}018 & \$68{,}335 & \$18{,}856 \\
2 & 435{,}449 & \$67{,}603 & \$3{,}470 & \$71{,}074 & \$38{,}468 \\
3 & 435{,}449 & \$69{,}970 & \$3{,}991 & \$73{,}960 & \$58{,}877 \\
4 & 435{,}449 & \$72{,}419 & \$4{,}589 & \$77{,}008 & \$80{,}126 \\
5 & 435{,}449 & \$74{,}953 & \$5{,}278 & \$80{,}231 & \$102{,}265 \\
\midrule
\textbf{Total} & \textbf{2{,}177{,}245} & \textbf{\$350{,}262} & \textbf{\$20{,}346} & \textbf{\$370{,}608} & \textbf{\$102{,}265} \\
\bottomrule
\end{tabular}
\end{table}

\noindent NPV of savings over 5 years (at 9.71\% discount rate): \textbf{\$81{,}161}. Adjusted payback: \textbf{1.31 years} (when evaluated against the lower-CAPEX demo scenario of \$25{,}000).

\subsection{CloudSim Rack-Level Analysis}

The CloudSim integration revealed critical airflow violations under AI workloads:

\begin{table}[h!]
\centering
\caption{Air-Side Economizer --- Rack-Level Hotspot Analysis (Mixed AI Workload)}
\label{tab:air_rack}
\renewcommand{\arraystretch}{1.3}
\begin{tabular}{lrrl}
\toprule
\textbf{Rack} & \textbf{Peak Load (kW)} & \textbf{Required CFM} & \textbf{Status} \\
\midrule
Rack 0 & 5.5 & 803 & \textcolor{red}{VIOLATION} (limit: 400 CFM) \\
Rack 1 & 5.4 & 792 & \textcolor{red}{VIOLATION} (limit: 400 CFM) \\
Rack 2 & 5.5 & 809 & \textcolor{red}{VIOLATION} (limit: 400 CFM) \\
Rack 3 & 10.0 & 1{,}464 & \textcolor{red}{HOTSPOT + VIOLATION} \\
Rack 4 & 10.0 & 1{,}456 & \textcolor{red}{HOTSPOT + VIOLATION} \\
\midrule
\textbf{Summary} & Avg: 7.05 kW & --- & 2/5 hotspots; CV = 0.30 \\
\bottomrule
\end{tabular}
\end{table}

\noindent All five racks exceed the 400 CFM per-rack airflow limit. Racks 3 and 4 are classified as hotspots (peak load 33\% above the 7.5 kW threshold). The system recommends upgrading to liquid cooling (direct-to-chip) for AI-density workloads.

\subsection{Hourly Operating Mode Distribution}

Based on the 8{,}760-hour simulation with real TMY3 weather data:

\begin{table}[h!]
\centering
\caption{Air-Side Economizer --- Annual Mode Distribution}
\label{tab:air_modes}
\renewcommand{\arraystretch}{1.3}
\begin{tabular}{lrrr}
\toprule
\textbf{Mode} & \textbf{Hours} & \textbf{\% of Year} & \textbf{Avg PUE} \\
\midrule
FULL\_ECON & $\sim$3{,}500 & $\sim$40\% & 1.23 \\
PARTIAL\_TRIM & $\sim$1{,}750 & $\sim$20\% & 1.33 \\
MECHANICAL\_ONLY & $\sim$3{,}510 & $\sim$40\% & 1.45 \\
\midrule
\textbf{Annual Average} & 8{,}760 & 100\% & \textbf{1.285} \\
\bottomrule
\end{tabular}
\end{table}

\subsection{Phase Gate Assessment}

\begin{table}[h!]
\centering
\caption{Air-Side Economizer --- Phase Gate Compliance}
\renewcommand{\arraystretch}{1.3}
\begin{tabular}{llp{7cm}}
\toprule
\textbf{Gate} & \textbf{Status} & \textbf{Notes} \\
\midrule
Thermal Compliance & \textcolor{orange}{MARGINAL} & Airflow violations at all racks under AI load \\
Water Constraint & \textcolor{green!60!black}{PASS} & Zero water consumption \\
Carbon Liability & \textcolor{red}{FAIL} & 219{,}921 kg CO$_2$/yr; high grid carbon intensity \\
Economic Viability & \textcolor{red}{FAIL} & Payback 491.9 yrs at full CAPEX (\$45k) \\
\bottomrule
\end{tabular}
\end{table}

\newpage

%% ============================================================
\section{Evaporative Cooling: Simulation Results}
%% ============================================================

\subsection{System Description}

The evaporative cooling system leverages the latent heat of vaporisation ($L_v = 2{,}260$ kJ/kg) to reduce supply air temperature. The simulation implements \textbf{Direct Evaporative Cooling (DEC)} and \textbf{Indirect Evaporative Cooling (IEC)} modes. Supply temperature is modelled as:
\[
T_{supply} = T_{db} - \eta_{actual} \times (T_{db} - T_{wb})
\]
where $\eta_{actual}$ is the saturation effectiveness (60--95\%). A lumped capacitance model prevents thermal spikes:
\[
C_{total}\frac{dT}{dt} = Q_{IT}(t) - Q_{cooling}(t)
\]
where $C_{total} = C_{rack} + C_{enclosure} = 15 + 30 = 45$ kJ/K.

\subsection{Annual Energy Results (8{,}760-Hour Simulation)}

\begin{table}[h!]
\centering
\caption{Evaporative Cooling --- Annual Energy \& Emissions Summary}
\label{tab:evap_annual}
\renewcommand{\arraystretch}{1.3}
\begin{tabular}{lrr}
\toprule
\textbf{Metric} & \textbf{Value} & \textbf{Unit} \\
\midrule
Total IT Energy & 141{,}173 & kWh \\
Fan Energy (Auxiliary) & 366 & kWh \\
DX Compressor Energy & 0 & kWh \\
Total Facility Energy & 141{,}539 & kWh \\
Average PUE & \textbf{1.002} & --- \\
Peak PUE & 1.004 & --- \\
Average CUE & 0.461 & kgCO$_2$/kWh$_{IT}$ \\
Total Carbon Emissions & 65{,}108 & kg CO$_2$ \\
Water Usage & \textbf{0} & Litres \\
Cooling Failure Hours & 0 & Hours \\
\bottomrule
\end{tabular}
\end{table}

\noindent The evaporative system achieves a near-ideal PUE of \textbf{1.002}, meaning only 0.2\% overhead above the IT load --- the best result of all three techniques. Zero DX compressor energy is consumed, and the system operates entirely on fan power.

\subsection{Economic Results}

\begin{table}[h!]
\centering
\caption{Evaporative Cooling --- Economic Analysis}
\label{tab:evap_econ}
\renewcommand{\arraystretch}{1.3}
\begin{tabular}{lrr}
\toprule
\textbf{Metric} & \textbf{Value} & \textbf{Unit} \\
\midrule
Annual Electricity Cost & \$20{,}523 & USD \\
Water Cost & \$0 & USD \\
Total Annual OpEx & \textbf{\$20{,}523} & USD \\
OpEx per kWh$_{IT}$ & \$0.1454 & USD/kWh \\
OpEx per Server (Annual) & \$213.78 & USD \\
Carbon Emissions & 65{,}108 & kg CO$_2$ \\
CO$_2$ per kWh$_{IT}$ & 0.461 & kg/kWh \\
\bottomrule
\end{tabular}
\end{table}

\noindent The evaporative system has the \textbf{lowest annual OpEx} of all three techniques at \$20{,}523/yr --- 79\% lower than Air-Side Economizer and 52\% lower than Chilled Water.

\subsection{Hourly Performance Profile}

The simulation runs in DEC/IEC mode throughout the year. Sample hourly data from the 8{,}760-hour run:

\begin{table}[h!]
\centering
\caption{Evaporative Cooling --- Sample Hourly Results (Hours 0--5)}
\label{tab:evap_hourly}
\renewcommand{\arraystretch}{1.3}
\begin{tabular}{rrrrrrl}
\toprule
\textbf{Hour} & \textbf{Amb. T (\textdegree C)} & \textbf{IT Load (kW)} & \textbf{Fan (kW)} & \textbf{Total (kW)} & \textbf{PUE} & \textbf{Mode} \\
\midrule
0 & 25.0 & 11.76 & 0.0142 & 11.774 & 1.0012 & DEC \\
1 & 25.0 & 11.76 & 0.0142 & 11.774 & 1.0012 & DEC \\
2 & 25.0 & 11.76 & 0.0142 & 11.774 & 1.0012 & DEC \\
3 & 25.0 & 11.76 & 0.0142 & 11.774 & 1.0012 & DEC \\
4 & 25.0 & 11.76 & 0.0142 & 11.774 & 1.0012 & DEC \\
5 & 25.0 & 11.798 & 0.0144 & 11.813 & 1.0012 & DEC \\
\bottomrule
\end{tabular}
\end{table}

\subsection{Cooling Assessment \& Thermal Compliance}

\begin{table}[h!]
\centering
\caption{Evaporative Cooling --- Cooling Assessment Checks}
\label{tab:evap_assess}
\renewcommand{\arraystretch}{1.3}
\begin{tabular}{llp{7cm}}
\toprule
\textbf{Check} & \textbf{Result} & \textbf{Detail} \\
\midrule
Heat Balance & \textcolor{red}{FAIL} & Cooling capacity 3.5\% below IT heat load at peak \\
Inlet Temperature & \textcolor{red}{FAIL} & Max inlet temp: 38.5\textdegree C (ASHRAE limit: 27\textdegree C) \\
Humidity & \textcolor{green!60!black}{PASS} & Supply humidity within acceptable range \\
Energy Efficiency & \textcolor{green!60!black}{PASS} & PUE 1.002 --- excellent \\
\midrule
\textbf{Overall Status} & \multicolumn{2}{l}{\textcolor{red}{\textbf{INSUFFICIENT\_COOLING}} (confidence: 95\%)} \\
\bottomrule
\end{tabular}
\end{table}

\noindent\textbf{Critical Finding:} Rack inlet temperatures reached \textbf{38.5\textdegree C}, exceeding the ASHRAE Class A1 limit of 27\textdegree C by 11.5\textdegree C. This triggers the thermal throttling model:
\[
\Phi_{adj} = \Phi_{base} \times (1 - \gamma \cdot \Delta T)
\]
At $\Delta T = 11.5$\textdegree C and $\gamma = 0.10$ (aggressive throttling above 29\textdegree C), computational throughput is reduced by up to \textbf{115\%} --- effectively a full shutdown. This makes standalone evaporative cooling \textbf{unsuitable for high-density AI workloads} without supplemental cooling.

\subsection{Engineering Recommendations}

\begin{enumerate}[noitemsep]
  \item Increase evaporative airflow by $\sim$5\% to close the 3.5\% capacity gap.
  \item Reduce IT load density or improve server efficiency.
  \item Deploy as a \textbf{pre-cooling stage} upstream of a chilled water system (hybrid configuration).
  \item Monitor dewpoint limits carefully --- direct evaporative cooling adds humidity to supply air.
\end{enumerate}

\subsection{Phase Gate Assessment}

\begin{table}[h!]
\centering
\caption{Evaporative Cooling --- Phase Gate Compliance}
\renewcommand{\arraystretch}{1.3}
\begin{tabular}{llp{7cm}}
\toprule
\textbf{Gate} & \textbf{Status} & \textbf{Notes} \\
\midrule
Thermal Compliance & \textcolor{red}{FAIL} & Inlet temp 38.5\textdegree C; 11.5\textdegree C above ASHRAE limit \\
Water Constraint & \textcolor{green!60!black}{PASS} & Zero water consumption in IEC mode \\
Carbon Liability & \textcolor{green!60!black}{PASS} & Lowest emissions: 65{,}108 kg CO$_2$/yr \\
Economic Viability & \textcolor{green!60!black}{PASS} & Lowest OpEx: \$20{,}523/yr \\
\bottomrule
\end{tabular}
\end{table}

\newpage

%% ============================================================
\section{Chilled Water System: Simulation Results}
%% ============================================================

\subsection{System Description}

The chilled water system uses the Energy Input Ratio (EIR) framework for plant operations. Efficiency is a function of chilled water supply temperature ($T_{chw}$) and condenser water temperature ($T_{cond}$). Variable-speed pumping and cooling tower fans follow affinity laws where power scales cubically with speed ($P \propto \text{speed}^3$). The system includes:

\begin{itemize}[noitemsep]
  \item Centrifugal chillers with EIR-based part-load performance curves
  \item Variable primary/secondary pumping loops
  \item Cooling tower with wet-bulb approach control
  \item CRAH units for in-row heat rejection
  \item Bayesian calibration for real-world performance matching
\end{itemize}

\subsection{Annual Energy Results (8{,}760-Hour Simulation)}

\begin{table}[h!]
\centering
\caption{Chilled Water System --- Annual Energy \& Emissions Summary}
\label{tab:chw_annual}
\renewcommand{\arraystretch}{1.3}
\begin{tabular}{lrr}
\toprule
\textbf{Metric} & \textbf{Value} & \textbf{Unit} \\
\midrule
Total Facility Energy & 288{,}475 & kWh \\
Cooling Load & 40{,}487 & kWh \\
Average PUE & 1.163 & --- \\
Water Usage & 515{,}331 & Litres \\
WUE & 1.786 & L/kWh$_{IT}$ \\
Average COP & 1.000 & --- \\
Peak Cooling Load & 5.853 & kW \\
Total Carbon Emissions & 158{,}196 & kg CO$_2$ \\
Annual OpEx & \$42{,}553 & USD \\
\bottomrule
\end{tabular}
\end{table}

\noindent\textbf{Note on COP:} The average COP of 1.0 reflects the simulation's conservative EIR baseline (COP$_{ref}$ = 3.0 at design conditions). The reported value represents the effective system COP including all auxiliary loads (pumps, cooling tower fans, CRAH units). Real-world chillers achieve COP 4--6 at part load.

\subsection{Economic Results}

\begin{table}[h!]
\centering
\caption{Chilled Water System --- Economic Analysis}
\label{tab:chw_econ}
\renewcommand{\arraystretch}{1.3}
\begin{tabular}{lrr}
\toprule
\textbf{Metric} & \textbf{Value} & \textbf{Unit} \\
\midrule
CAPEX & \$630{,}000 & USD \\
Annual OpEx & \$42{,}553 & USD \\
LCCP (15-year) & \$1{,}838{,}582 & USD \\
NPV & $-$\$1{,}092{,}355 & USD \\
Simple Payback Period & 50.0 & Years \\
\bottomrule
\end{tabular}
\end{table}

\noindent The negative NPV of $-$\$1{,}092{,}355 and 50-year payback reflect the high CAPEX (\$630{,}000) of a full chilled water plant. This is typical for enterprise-scale deployments where the investment is amortised over large IT loads (MW-scale). At the simulated scale ($\sim$20--25 kW IT load), the economics are unfavourable.

\subsection{Hourly Performance Profile}

\begin{table}[h!]
\centering
\caption{Chilled Water System --- Sample Hourly Results (Hours 0--4)}
\label{tab:chw_hourly}
\renewcommand{\arraystretch}{1.3}
\begin{tabular}{rrrrrrrr}
\toprule
\textbf{Hr} & \textbf{Amb. T} & \textbf{IT (kW)} & \textbf{Cool (kW)} & \textbf{Chiller (kW)} & \textbf{COP} & \textbf{Water (L)} & \textbf{Cost} \\
\midrule
0 & 24.0\textdegree C & 21.43 & 2.757 & 2.678 & 8.00 & 43.39 & \$2.64 \\
1 & 24.0\textdegree C & 21.43 & 2.757 & 2.678 & 8.00 & 43.39 & \$2.64 \\
2 & 24.0\textdegree C & 21.43 & 2.757 & 2.679 & 7.999 & 43.39 & \$2.64 \\
3 & 24.0\textdegree C & 21.43 & 2.758 & 2.679 & 7.998 & 43.39 & \$2.64 \\
4 & 24.0\textdegree C & 21.43 & 2.758 & 2.679 & 7.998 & 43.39 & \$2.64 \\
\bottomrule
\end{tabular}
\end{table}

\noindent The chilled water system maintains stable hourly COP around 8.0 at low ambient temperatures (24\textdegree C), demonstrating the efficiency advantage of vapour-compression refrigeration at near-design conditions.

\subsection{Water Consumption Analysis}

The chilled water system consumes \textbf{515{,}331 litres/year} (WUE = 1.786 L/kWh$_{IT}$). This is the highest water consumption of all three techniques and represents a significant operational constraint in water-scarce regions:

\begin{table}[h!]
\centering
\caption{Water Consumption Comparison by Region (Projected)}
\renewcommand{\arraystretch}{1.3}
\begin{tabular}{lrrl}
\toprule
\textbf{Region} & \textbf{Annual Water (L)} & \textbf{Water Cost} & \textbf{Scarcity Risk} \\
\midrule
Seattle (low scarcity) & 515{,}331 & Low & \textcolor{green!60!black}{Low} \\
Miami (moderate) & 515{,}331 & Moderate & \textcolor{orange}{Moderate} \\
Denver (moderate) & 515{,}331 & Moderate & \textcolor{orange}{Moderate} \\
Phoenix (high scarcity) & 515{,}331 & High + Penalty & \textcolor{red}{High} \\
\bottomrule
\end{tabular}
\end{table}

\subsection{Phase Gate Assessment}

\begin{table}[h!]
\centering
\caption{Chilled Water System --- Phase Gate Compliance (8{,}760-hr Run)}
\renewcommand{\arraystretch}{1.3}
\begin{tabular}{llp{7cm}}
\toprule
\textbf{Gate} & \textbf{Status} & \textbf{Notes} \\
\midrule
Thermal Compliance & \textcolor{green!60!black}{PASS} & Inlet temps maintained within ASHRAE A2 range \\
Water Constraint & \textcolor{green!60!black}{PASS} & 515{,}331 L/yr within site water budget \\
Carbon Liability & \textcolor{red}{FAIL} & 158{,}196 kg CO$_2$/yr exceeds carbon target \\
Economic Viability & \textcolor{red}{FAIL} & NPV $-$\$1.09M; payback 50 years \\
\bottomrule
\end{tabular}
\end{table}

\noindent\textbf{Key Advantage:} The chilled water system is the \textbf{only technique to achieve full thermal compliance} at ASHRAE Class A2 limits. It is the recommended choice for high-density AI workloads where server inlet temperature must be guaranteed below 27\textdegree C.

\newpage

%% ============================================================
\section{Comparative Analysis}
%% ============================================================

\subsection{Energy Efficiency Comparison}

\begin{table}[h!]
\centering
\caption{Energy Efficiency Metrics --- All Three Techniques}
\label{tab:compare_energy}
\renewcommand{\arraystretch}{1.4}
\begin{tabular}{lrrrl}
\toprule
\textbf{Metric} & \textbf{Air-Side Econ.} & \textbf{Evaporative} & \textbf{Chilled Water} & \textbf{Winner} \\
\midrule
Total Energy (kWh/yr) & 478{,}090 & \textbf{141{,}539} & 288{,}475 & Evaporative \\
Average PUE & 1.285 & \textbf{1.002} & 1.163 & Evaporative \\
Peak PUE & 1.45 & \textbf{1.004} & 1.20 & Evaporative \\
Cooling Overhead (\%) & 28.6\% & \textbf{0.2\%} & 16.3\% & Evaporative \\
Fan Energy (kWh/yr) & 105{,}965 & \textbf{366} & --- & Evaporative \\
Compressor Energy & High & \textbf{Zero} & Moderate & Evaporative \\
\bottomrule
\end{tabular}
\end{table}

\subsection{Environmental Impact Comparison}

\begin{table}[h!]
\centering
\caption{Environmental Metrics --- All Three Techniques}
\label{tab:compare_env}
\renewcommand{\arraystretch}{1.4}
\begin{tabular}{lrrrl}
\toprule
\textbf{Metric} & \textbf{Air-Side Econ.} & \textbf{Evaporative} & \textbf{Chilled Water} & \textbf{Winner} \\
\midrule
CO$_2$ Emissions (kg/yr) & 219{,}921 & \textbf{65{,}108} & 158{,}196 & Evaporative \\
CUE (kgCO$_2$/kWh$_{IT}$) & 0.591 & \textbf{0.461} & 0.549 & Evaporative \\
Water Usage (L/yr) & \textbf{0} & \textbf{0} & 515{,}331 & Air/Evap (tie) \\
WUE (L/kWh$_{IT}$) & \textbf{0} & \textbf{0} & 1.786 & Air/Evap (tie) \\
Carbon Tax/yr (\$) & \$27{,}710 & \textbf{\$8{,}270} & \$20{,}117 & Evaporative \\
\bottomrule
\end{tabular}
\end{table}

\subsection{Economic Comparison}

\begin{table}[h!]
\centering
\caption{Economic Metrics --- All Three Techniques}
\label{tab:compare_econ}
\renewcommand{\arraystretch}{1.4}
\begin{tabular}{lrrrl}
\toprule
\textbf{Metric} & \textbf{Air-Side Econ.} & \textbf{Evaporative} & \textbf{Chilled Water} & \textbf{Winner} \\
\midrule
CAPEX (USD) & \$45{,}000 & Low & \$630{,}000 & Evaporative \\
Annual OpEx (USD) & \$97{,}033 & \textbf{\$20{,}523} & \$42{,}553 & Evaporative \\
15-yr LCCP (USD) & --- & --- & \$1{,}838{,}582 & --- \\
NPV (USD) & \$81{,}161 & N/A & $-$\$1{,}092{,}355 & Air-Side \\
Payback (yrs) & 491.9 & N/A & 50.0 & --- \\
\bottomrule
\end{tabular}
\end{table}

\subsection{Thermal Compliance Comparison}

\begin{table}[h!]
\centering
\caption{Thermal Compliance --- All Three Techniques}
\label{tab:compare_thermal}
\renewcommand{\arraystretch}{1.4}
\begin{tabular}{lrrrl}
\toprule
\textbf{Metric} & \textbf{Air-Side Econ.} & \textbf{Evaporative} & \textbf{Chilled Water} & \textbf{Winner} \\
\midrule
Max Inlet Temp (\textdegree C) & $>$27 (violation) & 38.5 (violation) & \textbf{$\leq$27} & Chilled Water \\
ASHRAE A1 Compliance & \textcolor{orange}{Marginal} & \textcolor{red}{FAIL} & \textcolor{green!60!black}{PASS} & Chilled Water \\
Airflow Violations & All 5 racks & N/A & None & Chilled Water \\
Throttling Events & Yes & Yes & \textbf{None} & Chilled Water \\
Cooling Failure Hours & 0 & 0 & \textbf{0} & Tie \\
\bottomrule
\end{tabular}
\end{table}

\subsection{ML Recommender System Output}

The Random Forest recommender, trained on 8{,}760-hour simulation outputs, produces the following composite scores (higher = better overall fit):

\begin{table}[h!]
\centering
\caption{ML Recommender --- Composite Scores}
\label{tab:ml_scores}
\renewcommand{\arraystretch}{1.3}
\begin{tabular}{lrrrrl}
\toprule
\textbf{Technique} & \textbf{Score} & \textbf{Cost/yr} & \textbf{Emissions (kg)} & \textbf{Water (L)} & \textbf{Feasible} \\
\midrule
Evaporative & 0.200 & \$24{,}500 & 95{,}000 & 450{,}000 & Yes \\
Air-Side Econ. & 0.486 & \$26{,}000 & 100{,}000 & 5{,}000 & Yes \\
Chilled Water & \textbf{0.930} & \$27{,}235 & 102{,}084 & 294{,}088 & Yes \\
\bottomrule
\end{tabular}
\end{table}

\noindent\textbf{Interpretation:} The ML model assigns the highest composite score to \textbf{Chilled Water} (0.930) because it is the only technique that satisfies all thermal constraints --- the primary hard constraint for AI workloads. Evaporative cooling scores lowest (0.200) due to its thermal compliance failure, despite having the best energy and cost metrics.

\subsection{Radar Summary: Multi-Criteria Assessment}

\begin{table}[h!]
\centering
\caption{Multi-Criteria Assessment (1 = Best, 3 = Worst)}
\label{tab:radar}
\renewcommand{\arraystretch}{1.4}
\begin{tabular}{lrrr}
\toprule
\textbf{Criterion} & \textbf{Air-Side Econ.} & \textbf{Evaporative} & \textbf{Chilled Water} \\
\midrule
Energy Efficiency & 3 & \textbf{1} & 2 \\
Thermal Reliability & 2 & 3 & \textbf{1} \\
Carbon Emissions & 3 & \textbf{1} & 2 \\
Water Consumption & \textbf{1} & \textbf{1} & 3 \\
Capital Cost & 2 & \textbf{1} & 3 \\
Operating Cost & 3 & \textbf{1} & 2 \\
AI Workload Suitability & 2 & 3 & \textbf{1} \\
\midrule
\textbf{Total Score} & \textbf{16} & \textbf{11} & \textbf{14} \\
\textbf{Rank} & 3rd & \textbf{1st*} & 2nd \\
\bottomrule
\end{tabular}
\end{table}

\noindent *Evaporative ranks 1st on aggregate but fails the critical thermal compliance gate for AI workloads.

\newpage

%% ============================================================
\section{System Architecture Diagram}
%% ============================================================

\subsection{Overall Co-Simulation Architecture}

The following describes the full system architecture connecting all three cooling simulation engines with the shared CloudSim Plus workload layer and the ML recommender pipeline.

\begin{verbatim}
+------------------------------------------------------------------+
|              CLOUDSIM PLUS 8.x  (Workload Layer)                 |
|  Hosts: Physical servers with linear power models                |
|  VMs:   1:1 mapped, TimeShared scheduler                         |
|  Cloudlets: AI Training (1.8x) / Inference (1.4x) / Mixed (1.3x)|
|  Output: hourlyITLoadKW[8760]  (per-host, per-hour)              |
+------------------------------------------------------------------+
         |                    |                    |
         v                    v                    v
+----------------+  +------------------+  +--------------------+
| AIR-SIDE ECON  |  | EVAPORATIVE      |  | CHILLED WATER      |
| Spring Boot    |  | Python FastAPI   |  | Spring Boot        |
|                |  |                  |  |                    |
| AirEconomizer  |  | EvapCoolingModel |  | ChilledWaterPhysics|
| Model.java     |  | (DEC/IEC modes)  |  | .java              |
|                |  |                  |  |                    |
| Decision Tree: |  | Tsupply =        |  | EIR Framework:     |
| FULL_ECON      |  | Tdb - eta*(Tdb   |  | COP = f(Tchw,Tcond)|
| PARTIAL_TRIM   |  |   - Twb)         |  | Affinity Laws:     |
| MECH_ONLY      |  |                  |  | P ~ speed^3        |
|                |  | Lumped Cap:      |  |                    |
| Vreq = PIT*3160|  | C*dT/dt =        |  | Components:        |
| / (rho*Cp*dT)  |  | QIT - Qcooling   |  | Chiller + Pump     |
|                |  |                  |  | CoolingTower + CRAH|
+-------+--------+  +--------+---------+  +----------+---------+
        |                    |                       |
        v                    v                       v
+------------------------------------------------------------------+
|              RESULTS AGGREGATOR                                  |
|  Hourly: PUE, CUE, WUE, power breakdown, mode, violations       |
|  Annual: Energy, cost, emissions, water, phase gates             |
|  5-year: NPV, carbon tax escalation, climate scenarios           |
+------------------------------------------------------------------+
        |
        v
+------------------------------------------------------------------+
|              ML RECOMMENDER (Random Forest)                      |
|  Features: energy_kwh, water_L, cost_USD, emissions_kg,         |
|            violations, pue, wue, cue                             |
|  Output:   Ranked technique + composite score + why_recommended  |
+------------------------------------------------------------------+
        |
        v
+------------------------------------------------------------------+
|              REACT FRONTEND (Vite + TypeScript)                  |
|  Dashboard: Live KPIs, PUE/CUE charts, mode distribution        |
|  Simulation Detail: Hourly data tables, rack analysis            |
|  Reporting: PDF export, comparison tables                        |
+------------------------------------------------------------------+
\end{verbatim}

\subsection{Chilled Water System Internal Architecture}

\begin{verbatim}
+------------------------------------------------------------------+
|                  CHILLED WATER PLANT                             |
|                                                                  |
|  [Cooling Tower]                                                 |
|   Wet-bulb approach control                                      |
|   Fan power: P ~ speed^3 (affinity law)                          |
|   Evaporative heat rejection to atmosphere                       |
|        |                                                         |
|        v                                                         |
|  [Condenser Water Loop]  Tcond = f(Twb, approach)               |
|        |                                                         |
|        v                                                         |
|  [Centrifugal Chiller]                                           |
|   EIR = f(Tchw, Tcond) via bi-quadratic curves                   |
|   COP_ref = 3.0 (conservative baseline)                          |
|   Part-load: IPLV modelled via PLR curves                        |
|        |                                                         |
|        v                                                         |
|  [Chilled Water Loop]  Tchw = 7degC supply / 12degC return       |
|        |                                                         |
|        v                                                         |
|  [CRAH Units]  Computer Room Air Handlers                        |
|   Variable speed fans (affinity law)                             |
|   Sensible heat ratio control                                    |
|        |                                                         |
|        v                                                         |
|  [Server Racks]  IT Load = CloudSim hourlyITLoadKW               |
|   Thermal mass: Crack = 15 kJ/K                                  |
|   Inlet temp monitored vs ASHRAE 27degC limit                    |
+------------------------------------------------------------------+
\end{verbatim}

\subsection{Air-Side Economizer Control State Machine}

\begin{verbatim}
  START (each hour)
       |
       v
  Tout <= 24degC?
  /           \
YES            NO
 |              |
 v              v
RHout <= 60%?  MECHANICAL_ONLY
 /       \      Pcomp = QIT / COP(Tout)
YES      NO     Pfan  = Vreq * W/CFM
 |        |
 v        v
FULL_ECON  PARTIAL_TRIM
OA=100%    OA=20%
Pcomp=0    Pcomp = Qmech/COP
Pfan = Vreq * W/CFM * 1.15 (filter penalty)
       |
       v
  Airflow Violation Check:
  Vreq > Vmax?  --> VIOLATION: Recommend liquid cooling
       |
       v
  Compute: PUE, CUE, hourly cost, carbon
\end{verbatim}

\newpage

%% ============================================================
\section{Climate Resilience \& Sustainability Analysis}
%% ============================================================

\subsection{IPCC Climate Scenario Impact}

Performance is modelled against three IPCC RCP pathways for 2030. A 3\% cooling load increase per 1\textdegree C temperature rise is applied:

\begin{table}[h!]
\centering
\caption{Climate Scenario Impact on Annual Cooling Energy}
\label{tab:climate}
\renewcommand{\arraystretch}{1.3}
\begin{tabular}{lrrrr}
\toprule
\textbf{Scenario} & \textbf{Temp Offset} & \textbf{Air-Side (kWh)} & \textbf{Evap. (kWh)} & \textbf{Chilled W. (kWh)} \\
\midrule
RCP 2.6 (Baseline) & +0\textdegree C & 478{,}090 & 141{,}539 & 288{,}475 \\
RCP 4.5 (Moderate) & +1.5\textdegree C & 499{,}604 & 147{,}909 & 301{,}456 \\
RCP 8.5 (High) & +3.0\textdegree C & 521{,}118 & 154{,}278 & 314{,}438 \\
\bottomrule
\end{tabular}
\end{table}

\noindent Under RCP 8.5, the Air-Side Economizer is most severely impacted because higher ambient temperatures reduce free-cooling hours, forcing more mechanical operation. The Evaporative system is least sensitive to temperature rise (only 9\% increase vs 9\% for Chilled Water and 9\% for Air-Side).

\subsection{Carbon Tax Escalation (15\% Annual)}

\begin{table}[h!]
\centering
\caption{Carbon Tax Cost Projection (15\% Annual Escalation, IPCC 2030 Pathway)}
\label{tab:carbon_tax}
\renewcommand{\arraystretch}{1.3}
\begin{tabular}{lrrr}
\toprule
\textbf{Year} & \textbf{Air-Side Econ.} & \textbf{Evaporative} & \textbf{Chilled Water} \\
\midrule
Year 1 & \$27{,}710 & \$8{,}270 & \$20{,}117 \\
Year 2 & \$31{,}867 & \$9{,}511 & \$23{,}135 \\
Year 3 & \$36{,}647 & \$10{,}937 & \$26{,}605 \\
Year 4 & \$42{,}144 & \$12{,}578 & \$30{,}596 \\
Year 5 & \$48{,}466 & \$14{,}465 & \$35{,}185 \\
\midrule
\textbf{5-yr Total} & \textbf{\$186{,}834} & \textbf{\$55{,}761} & \textbf{\$135{,}638} \\
\bottomrule
\end{tabular}
\end{table}

\subsection{Water Stress Analysis}

\begin{table}[h!]
\centering
\caption{Water Stress Assessment by Technique and Region}
\label{tab:water_stress}
\renewcommand{\arraystretch}{1.3}
\begin{tabular}{lrrl}
\toprule
\textbf{Technique} & \textbf{Annual Water (L)} & \textbf{WUE} & \textbf{Phoenix Risk} \\
\midrule
Air-Side Economizer & 0 & 0.0 & \textcolor{green!60!black}{None} \\
Evaporative (IEC) & 0 & 0.0 & \textcolor{green!60!black}{None} \\
Chilled Water & 515{,}331 & 1.786 & \textcolor{red}{High --- Water Penalty Applied} \\
\bottomrule
\end{tabular}
\end{table}

\newpage

%% ============================================================
\section{Discussion \& Conclusions}
%% ============================================================

\subsection{Key Findings}

\begin{enumerate}
  \item \textbf{Evaporative Cooling achieves the best energy efficiency} (PUE 1.002) and lowest operating cost (\$20{,}523/yr) but \textbf{fails thermal compliance} for AI-density workloads. Rack inlet temperatures of 38.5\textdegree C exceed the ASHRAE Class A1 limit by 11.5\textdegree C, triggering aggressive CPU throttling.

  \item \textbf{Chilled Water is the only thermally compliant solution} for high-density AI workloads. It maintains inlet temperatures within ASHRAE A2 limits (18--27\textdegree C) at all hours. However, its high CAPEX (\$630{,}000) and water consumption (515{,}331 L/yr) are significant constraints.

  \item \textbf{Air-Side Economization is unsuitable as a standalone solution} for AI workloads. All five racks exhibit airflow violations (required CFM 2--4$\times$ the per-rack limit), and the system's high carbon emissions (219{,}921 kg CO$_2$/yr) result in a prohibitive carbon tax burden (\$27{,}710/yr).

  \item \textbf{The ML recommender correctly identifies Chilled Water} as the highest-scoring technique (0.930) when thermal compliance is weighted as a hard constraint, despite Evaporative having better energy and cost metrics.

  \item \textbf{Carbon tax escalation disproportionately penalises Air-Side Economization.} Over 5 years, its cumulative carbon tax (\$186{,}834) is 3.4$\times$ higher than Evaporative (\$55{,}761).
\end{enumerate}

\subsection{Recommended Hybrid Strategy}

Based on the simulation results, a \textbf{hybrid Evaporative + Chilled Water} configuration is recommended:

\begin{itemize}[noitemsep]
  \item Deploy evaporative pre-cooling to reduce condenser water temperature entering the chiller.
  \item Use chilled water as the primary cooling loop to guarantee ASHRAE compliance.
  \item This hybrid approach can achieve PUE $\approx$ 1.05--1.10 while maintaining full thermal compliance.
  \item Estimated water savings: 30--40\% vs. standalone chilled water.
\end{itemize}

\subsection{Limitations}

\begin{itemize}[noitemsep]
  \item The chilled water COP of 1.0 is conservative; real-world systems achieve COP 4--6 at part load.
  \item The evaporative simulation uses a constant ambient temperature (25\textdegree C) rather than TMY3 data, which may underestimate performance in cooler climates.
  \item Air-side economizer results use TMY3 weather data but a simplified constant COP = 3.0 for mechanical backup.
  \item All simulations assume a single-zone data centre; multi-zone hot/cold aisle containment would improve all results.
\end{itemize}

\subsection{Future Work}

\begin{enumerate}[noitemsep]
  \item Implement enthalpy-based economiser control (replacing dry-bulb/RH decision tree).
  \item Integrate real EIR part-load curves from manufacturer datasheets for chiller modelling.
  \item Extend the ML recommender to include hybrid configurations as candidate solutions.
  \item Validate simulation results against real data centre measurements (ASHRAE TC 9.9 benchmarks).
  \item Model multi-site deployment across Phoenix, Miami, Denver, and Seattle using full TMY3 datasets.
\end{enumerate}

\newpage

%% ============================================================
\section*{Appendix A: Governing Equations Reference}
%% ============================================================
\addcontentsline{toc}{section}{Appendix A: Governing Equations Reference}

\begin{table}[h!]
\centering
\caption{Complete Governing Equations Used Across All Three Simulations}
\renewcommand{\arraystretch}{1.5}
\begin{tabular}{p{4cm}p{5cm}p{4cm}}
\toprule
\textbf{Equation} & \textbf{Formula} & \textbf{Used In} \\
\midrule
Heat Balance & $C_{total}\frac{dT}{dt} = Q_{IT} - Q_{cooling}$ & All three \\
Airflow Requirement & $V_{req} = \frac{P_{IT,kW} \times 3160}{\rho \cdot C_p \cdot \Delta T}$ & Air-Side, Evap \\
DEC Supply Temp & $T_{supply} = T_{db} - \eta \cdot (T_{db} - T_{wb})$ & Evaporative \\
Fan Power & $P_{fan} = V_{CFM} \times \text{W/CFM} \times f_{filter}$ & Air-Side, Evap \\
Mechanical Cooling & $P_{mech} = Q_{mech} / \text{COP}$ & Air-Side, Chilled W \\
Affinity Law & $P \propto \text{speed}^3$ & Chilled Water \\
PUE & $\text{PUE} = P_{total} / P_{IT}$ & All three \\
CUE & $\text{CUE} = (P_{total} \times \text{CI}) / P_{IT}$ & All three \\
WUE & $\text{WUE} = \dot{m}_{water} / P_{IT}$ & Chilled Water \\
Throttling & $\Phi_{adj} = \Phi_{base}(1 - \gamma \cdot \Delta T)$ & All three \\
NPV & $\text{NPV} = \sum_{t=1}^{15} R_t / (1+i)^t$ & All three \\
Workload Model & $U(t) = [U_{diurnal}(t) \times M_{AI}] + \sigma(t)$ & All three \\
UPS Efficiency & $\eta = -0.05u^2 + 0.1u + 0.94$ & All three \\
\bottomrule
\end{tabular}
\end{table}

\vspace{1em}
\noindent\textbf{Constants:} $\rho = 1.2$ kg/m$^3$, $C_p = 1.006$ kJ/(kg$\cdot$K), $L_v = 2{,}260$ kJ/kg, $i = 0.0971$ (discount rate), $\gamma = 0.05$/\textdegree C (27--29\textdegree C), $\gamma = 0.10$/\textdegree C (29--32\textdegree C).

\end{document}
